% Arithmetic comparison. Independent whole-proof reading by /root completed;
% explicit positive-rotation marking added after that reader's requested check.
% No all-GT equality is asserted.

\subsection*{D23.F9. The exact scalar on the Galois subgroup}
Fix the usual ordered tangential braid base point
\[
 b_5=(0,t^4,t^3,t^2,t)
\]
and the monic hyperelliptic family
$Y^2=\prod_{j=1}^5(X-v_j)$ over the ordered configuration
space.  Choose the geometric path to the ordered CM
configuration $(1,\zeta_5,\zeta_5^2,\zeta_5^3,\zeta_5^4)$
that identifies its positive order-five rotation with the
standard class of $a=\sigma_1\sigma_2\sigma_3\sigma_4$.
The marking on the original character $\Omega_5$ is transported
along this path, with the order-five rotation represented by
$-q(a)$.  Then for $\sigma\in G_{\mathbb Q^{\rm ab}}$,
\[
 \mathcal C_5(\sigma)
       =(-1)^{\kappa_5(\sigma)}\Omega_5(\sigma).
 \tag{D23.F9.1}
\]
Consequently the character
\[
 \mathcal E_5(h)=(-1)^{\kappa_5(h)}
                  \mathcal C_5(h)\Omega_5(h)^{-1}
 \quad(h\in\mathcal H_5)
\]
is trivial on $G_{\mathbb Q^{\rm ab}}$.  This is a relative
character on the multiplier-one full GT group; no claim of its
triviality away from Galois or extension to the full gentle
group is made.

\begin{proof}
Let $M_\sigma$ be the arithmetic Tate action of the displayed
family at $b_5$. The exact base point and the arithmetic
Drinfeld action for every number of strands are recorded in
Nakamura, \emph{On arithmetic monodromy representations of
Eisenstein type in fundamental groups of once punctured
elliptic curves}, RIMS-1691 (2010), Section 5.4,
equation (5.4.2), printed p.~28. Arithmetic/topological
functoriality gives
\[
 q\Phi_\sigma=\operatorname{Ad}(M_\sigma)q.
\]
For multiplier one this is a symplectic action.  In addition
$M_\sigma e_1=e_1$ in the chain basis used by D23.T1.
To check the sign in this last assertion, let the first two
roots coalesce while the others remain at $t^3,t^2,t$.
The unit multiplying the double root at zero has leading
coefficient $-t^6$ and square root $it^3$.  Thus the node
splits over $\mathbb Q^{\rm ab}((t))$, with branches individually
fixed by $G_{\mathbb Q^{\rm ab}}$; its oriented vanishing
cycle is the first chain vector $e_1$.  The local annulus
has inertia $\widehat{\mathbb Z}(1)$, on which Galois acts
by the cyclotomic character, which is one here.  Hence the
cycle is fixed with its sign.  The uniqueness statement of
D23.T1 therefore gives $M_\sigma=V_\sigma$.

We specify the periodic-braid sign before comparing fibres.
Rotate the five branch points counterclockwise through
$2\pi/5$. Along the angular parameter $0\leq\theta\leq2\pi/5$
the isomorphisms from the initial curve to the rotating
family are
\[
 (X,Y)\longmapsto(e^{i\theta}X,e^{5i\theta/2}Y).
\]
Their endpoint is $(X,Y)\mapsto(\zeta_5X,-Y)$.
The positive periodic braid of this isotopy has the standard
class of $a$; choose the comparison path to identify these
positive classes. In the disc braid group both lifts have
exponent sum four, so a possible difference $z_5^e$ has
$20e=0$ and is trivial. The homological action of the
positive periodic braid is therefore exactly $q(a)$.
The order-five CM deck map $(X,Y)\mapsto(\zeta_5X,Y)$
differs from it by the hyperelliptic involution, whose
action on homology is $-I$. Its action is consequently
$-q(a)$, as required. This argument fixes the sign and
the direction of the cyclotomic marking, rather than only
the projective action.

Normalize the first two configuration coordinates to $0,1$.
As shown in D23.F1, its fundamental group is exactly
$\ker\ell_{12}\subset\widehat{PB}_5$.
At the base point the normalizing difference is $t^4$;
its fifth power has the fixed square root $t^{10}$, so this
normalization introduces no scalar on
$G_{\mathbb Q^{\rm ab}}$.

Let $p$ be the normalized CM point and transport its arithmetic
action to this normalized tangential base point along the fixed
path.  Its order-five automorphism is defined over
$\mathbb Q(\zeta_5)$ and is fixed by $\sigma$.
Consequently the path cocycle is a pure conjugator of $\bar a$.
By the pure centralizer calculation in D23.F1 it is precisely
$\alpha_5(\sigma)$.  Since the path lies in the normalized
configuration space, its lift in the pure braid group has
$\ell_{12}=0$ and is exactly $A_5(\sigma)$.
The standard base-point change formula for a local system now
identifies the Tate action at $p$, transported to the fixed
homology lattice, with
\[
 q(A_5(\sigma))^{-1}M_\sigma
       =\mathcal C_5(\sigma).
\]

Finally set $d=\zeta_5-1$.  The coordinate change
$X=1+dt$ in the original CM equation gives
\[
 X^5-1=d^5\prod_{j=0}^4
             \left(t-\frac{\zeta_5^j-1}{d}\right).
\]
The normalized monic family and the original curve are
therefore related by $Y=d^{5/2}Y_0$.
The only nontrivial scalar in their Tate actions is
$\sigma(\sqrt d)/\sqrt d$, because $d^2$ is already in
$\mathbb Q(\zeta_5)$.  D23.F4 identifies that scalar as
$(-1)^{\kappa_5(\sigma)}$.  The original Tate action is
$\Omega_5(\sigma)$ by D23.F3, giving (D23.F9.1).
All three factors are homomorphisms into the commutative
norm-one torus, so their product defines the stated relative
character.
\end{proof}

\paragraph{Scope of the comparison.}
The formula retains the specified tangential point, ordered
path, vanishing-cycle sign and base-point-change convention.
In particular the order-five deck rotation in $\Omega_5$
is identified with $-q(a)$. No equality between the three
characters on all of the full GT group is inferred from
their arithmetic equality.
